\documentclass[11pt]{article}
\usepackage[T1]{fontenc}
\usepackage[utf8]{inputenc}
\usepackage[english]{babel}
\usepackage{geometry}
\usepackage{enumitem}
\usepackage{amsmath,amssymb}
\usepackage{xcolor}
\usepackage{hyperref}

\geometry{margin=1in}
\emergencystretch=3em
\sloppy
\setlist[itemize]{leftmargin=*, itemsep=0.35em}
\setlist[enumerate]{leftmargin=*, itemsep=0.45em}
\newcommand{\loc}[1]{\texttt{\detokenize{#1}}}
\newcommand{\code}[1]{\texttt{\detokenize{#1}}}
\newcommand{\ac}[1]{#1}
\newcommand{\acp}[1]{#1s}

\title{Technical Review of the Manuscript}
\author{paper-review skill output}
\date{2026-06-17}

\begin{document}
\maketitle

\section*{Overall Assessment}

The manuscript is technically substantial and the main experimental narrative is coherent: Phase~1 establishes speed tracking, Phase~2 adds NOx and SOC objectives, and the WLTC experiment appropriately tempers the strongest environmental claims. The strongest parts are the explicit staged pipeline, the seed-wise reporting, and the unusually candid limitations section.

The main remaining risks are not small prose issues. They are implementation-facing ambiguities and claim-calibration problems that affect reproducibility and interpretation: the policy observes a deployment-unavailable thermal variable, SOC drift is reported with multiple meanings, shielded actions are not separated from learned policy behavior, and the WLTC conclusion relies on a selected seed under a comparison with different absolute initial SOC values. These should be fixed before final submission.

I explicitly checked the key equations, reward definitions, action mappings, appendix model interfaces, sign/units conventions, branch conventions, implementation formulas, result tables, figure captions, and bibliography hygiene. I found no inverse-trigonometric or coordinate-branch ambiguity; the main technical issues are instead notation drift, unit ambiguity, action-shield attribution, and overstated causal or benchmark claims.

\section*{Highest-Priority Fixes}

\begin{enumerate}
  \item Clarify that \(T_{\text{wall,SCR1}}\) is a privileged simulation observation, or add an ablation without it.
  \item Define SOC metrics once and use separate names for signed net change, absolute final drift, and peak excursion.
  \item Report action-shield intervention rates and distinguish raw actions from post-shield plant inputs in action-summary plots.
  \item Re-run or qualify the WLTC benchmark with matched initial SOC, or weaken the selected-seed superiority claim.
  \item Standardise fuel-command units and reward-weight notation, especially \(m_{fuel}\) and \(W_{\text{soc}^2}\).
\end{enumerate}

\section*{Major Technical and Factual Findings}

\begin{enumerate}
  \item \textbf{Privileged thermal observation weakens the deployment and generalisation claims.}
  \begin{itemize}
    \item \textbf{Location:} \loc{src/03_methodology.tex:77--86}, \loc{src/04_implementation.tex:154--168}, \loc{src/06_discussion.tex:63--66}.
    \item \textbf{Problem:} The methodology first says the observation is restricted to quantities that could plausibly be measured on a real vehicle, but the actual 12-dimensional observation includes \(T_{\text{wall,SCR1}}\). The discussion later states that this SCR1 wall-temperature variable is known not to be available during deployment.
    \item \textbf{Why it matters:} This is not just a limitation; it changes the formal problem. A non-recurrent policy with a privileged thermal proxy is not the same controller as a deployable non-recurrent policy using only real sensor signals. The WLTC ``zero-shot'' result may therefore show transfer across cycles in the surrogate setting, but not necessarily deployable generalisation.
    \item \textbf{Suggested fix:} State at the first observation-vector definition that \(T_{\text{wall,SCR1}}\) is a privileged simulation proxy. Either add an ablation without this variable, replace it with a measurable proxy, or revise abstract/conclusion language from deployment-oriented generalisation to ``surrogate-environment cycle transfer with one privileged thermal observation.''
  \end{itemize}

  \item \textbf{SOC terminology drifts across reward, objective, and results tables.}
  \begin{itemize}
    \item \textbf{Location:} \loc{src/03_methodology.tex:198--199}, \loc{src/04_implementation.tex:217}, \loc{src/04_implementation.tex:353--368}, \loc{src/05_results.tex:263--270}, \loc{src/05_results.tex:407--422}.
    \item \textbf{Problem:} The manuscript uses \(\Delta\text{SOC}_t=\text{SOC}_t-\text{SOC}_0\), signed \(\Delta\)SOC, \(|\Delta\text{SOC}|\), and \(\max_t|\Delta\mathrm{SOC}_t|\). These are all meaningful, but they are sometimes described generically as ``SOC drift'' or ``charge sustainment.'' Phase~1 reports signed mean \(\Delta\)SOC, Phase~2 reports both absolute final drift and peak drift, and the Optuna constraint uses peak drift.
    \item \textbf{Why it matters:} A controller can have small final \(|\Delta\text{SOC}|\) while making large transient SOC excursions, or a signed seed average can hide depletion and overcharging across seeds. This directly affects the claim that Phase~2 is charge-sustaining.
    \item \textbf{Suggested fix:} Introduce a short metric-definition table: signed net change \(\Delta\text{SOC}_{\mathrm{final}}\), absolute net drift \(|\Delta\text{SOC}_{\mathrm{final}}|\), and peak excursion \(\max_t|\Delta\text{SOC}_t|\). Use these names consistently in tables and captions, and specify which one defines ``charge-sustaining'' for each comparison.
  \end{itemize}

  \item \textbf{Fuel-command units and post-shield action signals are ambiguous.}
  \begin{itemize}
    \item \textbf{Location:} \loc{main.tex:208}, \loc{appendix.tex:15}, \loc{src/04_implementation.tex:66--75}, \loc{src/04_implementation.tex:154--186}, \loc{src/05_results.tex:300--303}.
    \item \textbf{Problem:} \(m_{fuel}\) is listed as mg/cyl/cycle in the symbols list and appendix, but several implementation tables use only mg. The Phase~2 PPO text also says the engine remains continuously on while the fuel-injection command spans roughly \(0\)--\(20\) mg, which conflicts with the stated engine-on valid range of \(3\)--\(70\) mg unless the plot shows raw pre-shield commands rather than plant inputs.
    \item \textbf{Why it matters:} The action mapping is a core reproducibility element. A reader implementing the environment needs to know whether the plotted values are raw agent actions, scaled commands, or post-shield plant inputs, and whether fuel mass is per cylinder per cycle or some other mass convention. This also affects whether reported operating points remain inside the LSTM validity domain.
    \item \textbf{Suggested fix:} Use one unit everywhere, preferably mg/cyl/cycle if that is the model input. Add a sentence defining the action-summary plots as pre-shield or post-shield. If the plotted fuel can reach zero while the engine is reported as on, explain the exception or correct the interpretation.
  \end{itemize}

  \item \textbf{The action shield and SOC clipping confound policy attribution.}
  \begin{itemize}
    \item \textbf{Location:} \loc{src/04_implementation.tex:262--269}, \loc{src/05_results.tex:297--303}, \loc{src/05_results.tex:367--373}, \loc{src/07_conclusion.tex:7--11}.
    \item \textbf{Problem:} The environment clips SOC, overrides EM2 torque near SOC boundaries, and enforces an engine-switch cooldown. The results and conclusion often describe the final behavior as learned charge sustainment or learned actuator strategy, but they do not report how often these shields intervene.
    \item \textbf{Why it matters:} If the shield frequently substitutes safe actions, the evaluated controller is a policy-plus-shield system, not the learned policy alone. This is acceptable, but the contribution and comparison to the rule-based baseline must be described accordingly.
    \item \textbf{Suggested fix:} Report per-cycle intervention counts: SOC clips, EM2 torque overrides, rejected engine switches, and forced \((n_{ICE},m_{fuel})\) defaults. Add either a no-shield diagnostic evaluation or explicitly rename the reported controller as a shielded learned controller.
  \end{itemize}

  \item \textbf{Reward clipping can create a mismatch with the cumulative NOx evaluation metric.}
  \begin{itemize}
    \item \textbf{Location:} \loc{src/04_implementation.tex:209--237}, especially \(\tilde e_{\mathrm{NO}_{x},t}=\min(\dot m_{\mathrm{NO}_{x},\mathrm{tp},t}/c_{\mathrm{NO}_{x},\mathrm{rate}},1)\).
    \item \textbf{Problem:} The reward caps the per-step NOx penalty above \(0.4\) g/s, while evaluation uses uncapped cumulative NOx mass in \loc{src/04_implementation.tex:360--365}.
    \item \textbf{Why it matters:} If high surrogate NOx spikes occur, the policy receives no additional penalty above the cap even though the final metric still accumulates the full mass. This can create a reward--metric mismatch and may hide rare but important emission bursts.
    \item \textbf{Suggested fix:} Report the fraction of evaluation steps where the NOx term is clipped for the selected policies. If clipping is common, consider a smooth robust transform that still distinguishes severe spikes, or align the reward more closely with the uncapped cumulative metric.
  \end{itemize}

  \item \textbf{The operating-point explanation overstates causality.}
  \begin{itemize}
    \item \textbf{Location:} \loc{src/05_results.tex:422}, \loc{src/05_results.tex:452--455}, \loc{src/06_discussion.tex:27--30}.
    \item \textbf{Problem:} The manuscript states that NOx rate rises monotonically with torque and that Phase~1 policies burned extra fuel at high engine load ``purely'' to overcharge the battery. The cited evidence is a map overlay and a Pearson correlation of 0.84, which supports a strong association but not monotonicity or a complete causal decomposition.
    \item \textbf{Why it matters:} This mechanism is central to the claim that Phase~2 reduces fuel and NOx together. Overstating it makes the result look more explanatory than the current evidence supports, especially because NOx depends on speed, fuel, aftertreatment state, temperature, and transient history, not torque alone.
    \item \textbf{Suggested fix:} Replace causal wording with ``is consistent with'' unless additional evidence is added. A stronger version would include an energy-balance check, intervention analysis, or multivariate map/partial-dependence analysis over torque, speed, fuel, and SCR temperature.
  \end{itemize}

  \item \textbf{The WLTC benchmark mixes different absolute initial SOC values and then relies on selected-seed language.}
  \begin{itemize}
    \item \textbf{Location:} \loc{src/05_results.tex:494--507}, \loc{src/05_results.tex:526--572}, \loc{main.tex:109}, \loc{src/07_conclusion.tex:11}.
    \item \textbf{Problem:} The RL policies are compared at \(\mathrm{SOC}\approx0.70\), while the rule-based isoSOC run is at \(\mathrm{SOC}\approx0.56\). The manuscript correctly says the 10-seed averages do not beat the rule-based controller environmentally, but the abstract and conclusion emphasize that SAC seed~7 reaches 80 mg/km and improves speed tracking.
    \item \textbf{Why it matters:} Charge-neutrality is not identical to equal initial SOC if the controller, battery limits, engine-off availability, or EM2 usage depend on absolute SOC. Also, a selected-seed result is weaker evidence than a robust population-level result.
    \item \textbf{Suggested fix:} Evaluate the rule-based and RL controllers from matched initial SOC values, or explain why absolute SOC is irrelevant for this plant and controller. In the abstract and conclusion, phrase the result as ``one selected SAC seed'' rather than as a general agent-level environmental improvement.
  \end{itemize}

  \item \textbf{Reward-weight notation drifts in the Phase~2 results table.}
  \begin{itemize}
    \item \textbf{Location:} \loc{src/04_implementation.tex:212}, \loc{src/04_implementation.tex:249--257}, \loc{src/05_results.tex:232--245}, \loc{src/06_discussion.tex:142}.
    \item \textbf{Problem:} The reward and implementation define \(W_{\text{soc}^2}\), but the best Phase~2 Optuna-trials table in the manuscript labels the column \(W_{\text{soc}}\).
    \item \textbf{Why it matters:} The column header suggests a linear SOC-deviation weight, while the actual reward weight multiplies \((\Delta\text{SOC}_t)^2\). This is a notation-drift issue in an implementation-facing table.
    \item \textbf{Suggested fix:} Change the table header to \(W_{\text{soc}^2}\) or rename the parameter everywhere to a clearer form such as \(W_{\Delta\mathrm{SOC}^2}\).
  \end{itemize}

  \item \textbf{The composite score \(J\) should be described explicitly as a unit-weighted ranking score.}
  \begin{itemize}
    \item \textbf{Location:} \loc{src/04_implementation.tex:350--379}, \loc{src/05_results.tex:300--301}, \loc{src/05_results.tex:370--371}, \loc{src/05_results.tex:552--553}.
    \item \textbf{Problem:} \(J\) adds grams of NOx, squared km/h violations, and squared SOC-fraction violations after applying heuristic coefficients. This is fine for optimisation, but the manuscript sometimes presents scores such as \(J=4.23\) as if the value were directly interpretable.
    \item \textbf{Why it matters:} The numerical value of \(J\) depends on the chosen tolerances, units, and coefficients; it is not a physical metric and may not transfer unchanged from the staircase cycle to WLTC selection.
    \item \textbf{Suggested fix:} State that \(J\) is a dimensionally weighted ranking score only. If it is also used for WLTC seed selection, specify the exact WLTC version of the score, including distance normalisation, thresholds, and whether \(\max|\Delta\text{SOC}|\) or final drift is used.
  \end{itemize}
\end{enumerate}

\section*{Meaningful Editorial and Claim-Calibrating Findings}

\begin{enumerate}
  \item \textbf{Conclusion overstates the WLTC result.}
  \begin{itemize}
    \item \textbf{Location:} \loc{src/07_conclusion.tex:11}.
    \item \textbf{Problem:} ``The zero-shot WLTC comparison further proved the agent's ability to generalise'' is too strong because the population-level NOx/fuel outcome is unfavourable and only one selected seed reaches the highlighted NOx result.
    \item \textbf{Why it matters:} This is the final claim readers will remember, so it should match the cautious discussion.
    \item \textbf{Suggested fix:} Rewrite as: ``The zero-shot WLTC comparison provided evidence that the learned policies transfer their speed-tracking behaviour to an unseen driving cycle, but it did not demonstrate a population-level environmental advantage over the rule-based controller. One selected SAC seed reached 80 mg/km NOx while remaining charge-sustaining, at higher fuel consumption.''
  \end{itemize}

  \item \textbf{The final WLTC-methodology sentence contains several errors and obscures the point.}
  \begin{itemize}
    \item \textbf{Location:} \loc{src/03_methodology.tex:214--217}.
    \item \textbf{Problem:} The sentence contains ``bridge the gap simulation to reality gap,'' ``sitation,'' ``havent,'' and ``effectivity.''
    \item \textbf{Why it matters:} This is the methodological justification for the WLTC experiment, so the wording should be precise.
    \item \textbf{Suggested fix:} ``This final comparison reduces the simulation-to-reality gap by assessing the learned controllers on a more realistic driving cycle rather than only on synthetic staircase profiles. It evaluates whether the final models generalise to situations not seen during training and compares their speed tracking, NOx reduction, and SOC sustainment with a rule-based solution developed by Frey et al.''
  \end{itemize}

  \item \textbf{Some sustainability wording conflicts with the WLTC fuel result.}
  \begin{itemize}
    \item \textbf{Location:} \loc{src/06_discussion.tex:150--160}.
    \item \textbf{Problem:} The SDG~12 bullet says the work makes more economical use of fuel per kilometre travelled, while the WLTC population-level comparison reports 1.6--2.2\(\times\) higher fuel consumption than the rule-based baseline.
    \item \textbf{Why it matters:} This can read as greenwashing despite the manuscript's otherwise careful caveats.
    \item \textbf{Suggested fix:} Restrict the statement to the controlled staircase Phase~2 comparison, or rewrite as a future objective: ``A future multi-objective formulation including fuel or CO2 could support SDG~12 if it reduces fuel use relative to established baselines.''
  \end{itemize}

  \item \textbf{The limitations subsection has a malformed lead sentence.}
  \begin{itemize}
    \item \textbf{Location:} \loc{src/06_discussion.tex:79--86}.
    \item \textbf{Problem:} ``Furthermore than the general Phase~2 transfer trategy'' contains both an idiom error and a typo.
    \item \textbf{Why it matters:} The paragraph introduces important methodological limitations and should not look draft-like.
    \item \textbf{Suggested fix:} ``Beyond the general Phase~2 transfer strategy, two specific adjustments could have improved the overall quality of the work.''
  \end{itemize}
\end{enumerate}

\section*{Proofing-Sweep Items}

The bundled proofing scan was run on the manuscript PDF and on comment-stripped source text. The PDF hits were table-of-contents dot leaders; the source hits were ADS bibcode/URL strings and acceptable semicolon cases. After spot-checking, I found no actionable duplicate-comma, malformed-quote, product-capitalisation, or \code{arctan(x/y)} branch-pattern hits. The following high-confidence copyedits remain worth fixing:

\begin{itemize}
  \item \textbf{Location:} \loc{src/01_introduction.tex:5}. \textbf{Problem:} Missing space in ``activity,measured.'' \textbf{Fix:} ``activity, measured.''
  \item \textbf{Location:} \loc{src/01_introduction.tex:23}. \textbf{Problem:} ``unfeasable.'' \textbf{Fix:} ``unfeasible'' or ``infeasible.''
  \item \textbf{Location:} \loc{src/02_theory.tex:30}. \textbf{Problem:} duplicated word in ``enables enables.'' \textbf{Fix:} remove one ``enables.''
  \item \textbf{Location:} \loc{src/04_implementation.tex:102}. \textbf{Problem:} ``intitialisation.'' \textbf{Fix:} ``initialisation.''
  \item \textbf{Location:} \loc{src/05_results.tex:426}. \textbf{Problem:} ``TD3 plateaus at ($10$\,g).'' \textbf{Fix:} ``TD3 plateaus at approximately $10$\,g.''
  \item \textbf{Location:} \loc{src/06_discussion.tex:9}. \textbf{Problem:} ``the central message is, that,''. \textbf{Fix:} ``the central message is that''.
  \item \textbf{Location:} \loc{src/06_discussion.tex:50}. \textbf{Problem:} ``The WLTC result therefore demonstrate'' and ``persiting.'' \textbf{Fix:} ``The WLTC results therefore demonstrate'' and ``persisting.''
\end{itemize}

\section*{Items Needing External Verification or Clearer Source Framing}

\begin{itemize}
  \item \textbf{Current transport and market statistics.} The 2023 EU GHG share and early-2026 combustion-vehicle market-share statements in \loc{src/01_introduction.tex:3--5} rely on time-sensitive sources. Add access dates and ensure the cited values match the exact geography and vehicle category.
  \item \textbf{Live carbon-intensity values.} The Spain/Germany electricity-carbon-intensity comparison in \loc{src/06_discussion.tex:104} uses a live source. Add the exact access date/time and clarify that the values are snapshots, not annual averages.
  \item \textbf{Regulatory framing.} The Euro~6/Euro~7 language in \loc{src/01_introduction.tex:43}, \loc{src/05_results.tex:426--431}, and \loc{src/05_results.tex:572} should distinguish an indicative mg/km comparison from regulatory compliance, especially for the synthetic staircase cycle and a surrogate-only evaluation.
  \item \textbf{Rule-based benchmark conditions.} The WLTC comparison should state precisely whether Frey et al.'s rule-based result is from the same surrogate stack, a GT-Suite simulation, or another setup, and whether the same initial conditions and actuator constraints apply.
\end{itemize}

\end{document}